\appendix
\chapter{Curriculum and Course Coverage Map}
\label{ch:course-coverage-map}

This appendix makes the curriculum auditable.
It maps the supplied external machine-learning course outcomes to the AI Almanach
and adds the broader competencies required to build dependable AI systems.

\begin{table}[htbp]
\centering
\caption{Coverage of the supplied external machine-learning course.}
\label{tab:external-course-coverage}
\begin{tblr}{
    width=\textwidth,
    colspec={X[l] Q[l,0.27\textwidth] X[l]},
    row{1}={font=\bfseries, bg=\maincolor!15!white},
    hlines,
    vlines}
Course requirement & Primary location & Evidence of mastery \\
Statistics, derivatives, and matrices for ML &
Mathematical Foundations &
Derive a gradient, check tensor dimensions, and quantify uncertainty \\
Classification and regression &
Machine Learning &
Implement baselines and compare metrics on held-out data \\
Model optimisation &
Deep Learning; Evaluation and MLOps &
Tune only on validation data and report the search protocol \\
Decision trees and forests &
Machine Learning &
Explain a split, control complexity, and measure generalisation \\
Neural networks and support vector machines &
Machine Learning; Deep Learning &
Calculate a toy forward pass and diagnose optimisation behaviour \\
Clustering and recommender systems &
Machine Learning; E-Commerce and Marketing &
Evaluate without pretending that unlabeled structure has one true answer \\
Develop, implement, and scientifically evaluate an ML approach &
Evaluation and MLOps; Practical LLM Engineering &
Reproduce a baseline-to-candidate experiment with uncertainty and error
analysis \\
\end{tblr}
\end{table}

\section{Competencies added by the AI Almanach}

\begin{itemize}
    \item formulate a decision problem before selecting a model;
    \item create data sheets, model cards, and reproducible experiment records;
    \item detect leakage, distribution shift, shortcut learning, and metric
          mismatch;
    \item measure latency, throughput, memory, energy, and operational
          reliability;
    \item evaluate generative outputs with task-specific, human, and safety
          tests;
    \item reason about privacy, security, functional safety, fairness, and law;
    \item communicate evidence and uncertainty to technical and non-technical
          readers.
\end{itemize}

\section{Alignment with current curriculum frameworks}

The curriculum also uses the ACM/IEEE/AAAI computing competencies as a
technical breadth check and UNESCO's AI competency framework as a human-centred
literacy check\textsuperscript{\cite{acmCS2023Curricula2024,unescoAICompetencyFramework2024}}.
These frameworks are complementary rather than interchangeable: one helps
audit computing depth, while the other foregrounds human agency, ethics, AI
techniques, and system design.

\begin{table}[htbp]
\centering
\caption{Framework alignment used to audit the AI Almanach curriculum.}
\label{tab:curriculum-framework-alignment}
\begin{tblr}{
    width=\textwidth,
    colspec={Q[l,0.22\textwidth] X[l] X[l]},
    row{1}={font=\bfseries, bg=\maincolor!15!white},
    hlines,
    vlines}
Audit lens & Almanach coverage & Observable evidence \\
Technical foundations &
Mathematics, programming, data, ML, deep learning, RL, vision, and NLP &
Trace a mechanism, implement a baseline, and test generalisation \\
Systems and operations &
Evaluation, MLOps, LLM engineering, security, and industrial systems &
Specify release gates, operational budgets, monitoring, and rollback \\
Human agency and ethics &
Trustworthy AI, healthcare, finance, commerce, law, and visualization &
Identify affected actors, contestability, harms, uncertainty, and duties \\
Design and creation &
Domain chapters, multi-agent systems, venture building, and scientific
writing &
Produce a bounded solution, evidence record, and reproducible argument \\
\end{tblr}
\end{table}

\section{Definition of done for a chapter}

A chapter is considered reviewed only when its important concepts have
prerequisites, definitions, dimensional checks, at least one worked example,
failure modes, evaluation guidance, and traceable sources.
The chapter must compile without hidden local assets and must not contain
unresolved placeholders.
Time-sensitive claims must include a verification date or be rewritten as
stable principles.
